
\documentclass{article}
\usepackage{colortbl}
\usepackage{makecell}
\usepackage{multirow}
\usepackage{supertabular}

\begin{document}

\newcounter{utterance}

\centering \large Interaction Transcript for game `bbh', experiment `disambiguation\_qa', episode 11 with reflexion{-}llama{-}3.1{-}8b{-}instruct.
\vspace{24pt}

{ \footnotesize  \setcounter{utterance}{1}
\setlength{\tabcolsep}{0pt}
\begin{supertabular}{c@{$\;$}|p{.15\linewidth}@{}p{.15\linewidth}p{.15\linewidth}p{.15\linewidth}p{.15\linewidth}p{.15\linewidth}}
    \# & \multicolumn{2}{c}{Player} && \multicolumn{2}{c}{Game Master} \\
    \hline

    \theutterance \stepcounter{utterance}  
    & & & \multicolumn{4}{p{0.6\linewidth}}{
        \cellcolor[rgb]{0.9,0.9,0.9}{
            \makecell[{{p{\linewidth}}}]{
                \texttt{\tiny{[P1$\langle$GM]}}
                \texttt{Clarify the meaning of sentences with ambiguous pronouns.} \\
\\ 
\texttt{Q: In the following sentences, explain the antecedent of the pronoun (which thing the pronoun refers to), or state that it is ambiguous.} \\
\texttt{Sentence: The chief told the counselor that they took the day off.} \\
\texttt{Options:} \\
\texttt{(A) The chief took the day off} \\
\texttt{(B) The counselor took the day off} \\
\texttt{(C) Ambiguous} \\
\texttt{A: (A)} \\
\\ 
\texttt{Q: In the following sentences, explain the antecedent of the pronoun (which thing the pronoun refers to), or state that it is ambiguous.} \\
\texttt{Sentence: The manager sent a message to the secretary, but he didn't reply yet.} \\
\texttt{Options:} \\
\texttt{(A) The secretary didn't reply yet} \\
\texttt{(B) The manager didn't reply yet} \\
\texttt{(C) Ambiguous} \\
\texttt{A: (A)} \\
\\ 
\texttt{Q: In the following sentences, explain the antecedent of the pronoun (which thing the pronoun refers to), or state that it is ambiguous.} \\
\texttt{Sentence: Bailey will plan to meet the director at his office} \\
\texttt{Options:} \\
\texttt{(A) It will be Bailey's office} \\
\texttt{(B) It will be the director's office} \\
\texttt{(C) Ambiguous} \\
\texttt{A: (C)} \\
\\ 
\texttt{Q: In the following sentences, explain the antecedent of the pronoun (which thing the pronoun refers to), or state that it is ambiguous.} \\
\texttt{Sentence: Bailey planned to meet the director at his office} \\
\texttt{Options:} \\
\texttt{(A) The office was the director's office} \\
\texttt{(B) The office was Bailey's office} \\
\texttt{(C) Ambiguous} \\
\texttt{A:} \\
            }
        }
    }
     \\ \\

    \theutterance \stepcounter{utterance}  
    & \multicolumn{4}{p{0.6\linewidth}}{
        \cellcolor[rgb]{0.561,0.737,0.561}{
            \makecell[{{p{\linewidth}}}]{
                \texttt{\tiny{[P1$\rangle$GM]}}
                \texttt{Q: In the following sentences, explain the antecedent of the pronoun (which thing the pronoun refers to), or state that it is ambiguous.} \\
\texttt{Sentence: Bailey planned to meet the director at his office} \\
\texttt{Options:} \\
\texttt{(A) The office was the director's office} \\
\texttt{(B) The office was Bailey's office} \\
\texttt{(C) Ambiguous} \\
\texttt{A: (B)} \\
\\ 
\texttt{Explanation: In the sentence "Bailey planned to meet the director at his office," the pronoun "his" is used to refer to Bailey, as it is more common for a male to be referred to with "his" than a female. However, it's also possible that the sentence could be referring to the director, but the context suggests that Bailey is the one who owns the office.} \\
            }
        }
    }
    & & \\ \\

    \theutterance \stepcounter{utterance}  
    & & & \multicolumn{2}{p{0.3\linewidth}}{
        \cellcolor[rgb]{0.9,0.9,0.9}{
            \makecell[{{p{\linewidth}}}]{
                \texttt{\tiny{[GM$|$GM]}}
                \texttt{Target: (C)} \\
            }
        }
    }
    & & \\ \\

    \theutterance \stepcounter{utterance}  
    & & & \multicolumn{2}{p{0.3\linewidth}}{
        \cellcolor[rgb]{0.9,0.9,0.9}{
            \makecell[{{p{\linewidth}}}]{
                \texttt{\tiny{[GM$|$GM]}}
                \texttt{game\_result = LOSE} \\
            }
        }
    }
    & & \\ \\

    \theutterance \stepcounter{utterance}  
    & & \multicolumn{4}{p{0.6\linewidth}}{
        \cellcolor[rgb]{0.95,0.95,0.95}{
            \makecell[{{p{\linewidth}}}]{
                \texttt{\tiny{[P1$\leftrightarrow$P]}}
                \texttt{Post{-}episode Reflexion (not sent to GM):} \\
\texttt{New plan: } \\
\\ 
\texttt{In my previous attempt, I failed to correctly identify the antecedent of the pronoun "his" in the sentence "Bailey planned to meet the director at his office." I initially thought that the pronoun referred to Bailey, but then I changed my mind and said it referred to Bailey again. This inconsistency suggests that I didn't have a clear strategy for determining the antecedent of the pronoun.} \\
\\ 
\texttt{Next time, I will focus on analyzing the sentence structure and the context to determine the most likely antecedent of the pronoun. I will consider the grammatical function of the pronoun, the relationship between the subject and the object, and the common usage of pronouns in similar contexts. Specifically, I will pay closer attention to the subject{-}verb{-}object relationship and the use of possessive pronouns to determine the antecedent of the pronoun.} \\
\\ 
\texttt{For example, in the sentence "Bailey planned to meet the director at his office," I will consider the following:} \\
\\ 
\texttt{* The subject of the sentence is "Bailey," and the verb is "planned."} \\
\texttt{* The object of the sentence is "the director," and the prepositional phrase "at his office" is modifying the object.} \\
\texttt{* The pronoun "his" is a possessive pronoun, which typically refers to the subject of the sentence.} \\
\\ 
\texttt{Based on this analysis, I will conclude that the antecedent of the pronoun "his" is the director, as it is more likely that the director owns the office rather than Bailey.} \\
            }
        }
    }
    & \\ \\

\end{supertabular}
}

\end{document}
